\emph{Preliminaries.} By Theorem \claimref{synthesis.signed-directions-and-the-integer-clifford-algebra} the integer algebra is free on the blades, each blade is invertible, and $\dagger$, the linear extension of blade inversion, is an anti-automorphism; $\langle xy^{\dagger}\rangle_{0}$ is the sum of products of coefficients, so $\langle x^{\dagger}\rangle_{0} = \langle x\rangle_{0}$, $\langle ab\rangle_{0} = \langle ba\rangle_{0}$, a self-adjoint blade has square $1$, and the scalar part of a product of two blades is zero unless they are equal up to sign. With $S = B_{1} + B_{2}$, $D = B_{1} - B_{2}$: $C = A_{1}S + A_{2}D$, $C^{2} = S^{2} + D^{2} + A_{1}A_{2}SD + A_{2}A_{1}DS = 4 - [A_{1},A_{2}][B_{1},B_{2}]$. By clause (ii) of Theorem \claimref{synthesis.mode-split-of-the-clifford-algebra-and-its-klein-twist-to-the-tensor-product}, $e_{2}p_{1}$ has square $1$ and commutes with $e_{1}$ and $p_{1}$; it is self-adjoint and anticommutes with $p_{2}$; $p_{1}$ and $e_{1}$ anticommute. So $[A_{1},A_{2}] = 2A_{1}A_{2}$, $[B_{1},B_{2}] = 2B_{1}B_{2}$, $C^{2} = 4 + 4K$; $K^{\dagger} = -B_{2}B_{1}A_{2}A_{1} = K$ and $K^{2} = KK^{\dagger} = 1$.

\emph{(i).} The four products are the blades $p_{1}p_{2}$, $p_{1}e_{2}p_{1} = e_{2}$, $e_{1}p_{2}$ and $e_{1}e_{2}p_{1}$, up to sign: they are distinct, none is a scalar, each is self-adjoint of square $1$, and none equals $\pm$ a single $A_{i}$ or $B_{j}$ or $\pm K$ (the check lists them). $1 + C$ is self-adjoint with integer coefficients, so $\lambda$ is integer-valued on the integer algebra and $\lambda(x^{\dagger}) = \langle(x(1 + C))^{\dagger}\rangle_{0} = \lambda(x)$. For a blade $x$ of square $1$, $\lambda(x)$ is the coefficient of $x$ in $1 + C$: this gives $\lambda(1) = 1$, the four values $1, 1, 1, -1$, and $0$ on the $A_{i}$, the $B_{j}$ and $K$. Expanding, $\lambda(\tfrac14(1 + aA_{i})(1 + bB_{j})) = \tfrac14(1 + ab\,\lambda(A_{i}B_{j}))$, which is $\tfrac12$ or $0$; summing over $b$ gives $\tfrac12$, independent of $j$, and likewise over $a$. Finally $\lambda((C - t)(C - t)^{\dagger}) = \lambda(C^{2}) - 2t\lambda(C) + t^{2} = 4 + 4\lambda(K) - 8t + t^{2} = t^{2} - 8t + 4$, negative for $4 - 2\sqrt{3} < t < 4 + 2\sqrt{3}$, in particular at the integers $1$ to $7$.

\emph{(ii).} $\varphi((C - t)^{2}) = \varphi(C^{2}) - 2t\varphi(C) + t^{2} \ge 0$ at $t = \varphi(C)$ gives $\varphi(C)^{2} \le \varphi(C^{2})$. $(1 - K)^{2} = 2 - 2K$, so $\varphi((1 - K)^{2}) \ge 0$ gives $\varphi(K) \le 1$ and $\varphi(C^{2}) = 4 + 4\varphi(K) \le 8$. $C - t$ contains all four products $A_{i}B_{j}$; $K$ is, up to sign, $A_{1}B_{1}\cdot A_{2}B_{2}$, a product of observables of two setting pairs. For $\lambda$: $\lambda((1 - K)^{2}) = 2 - 2\lambda(K) = 2$. Since $\lambda$ is a normalised linear functional on the algebra, satisfies the identity (which is an identity of elements) and has $\lambda(C)^{2} = 16 > 8$, no argument from the algebra and linearity alone gives the bound.

\emph{(iii).} $p_{3}$ is self-adjoint of square $1$ and commutes with the generators of modes $1$ and $2$, so $P_{\pm}$ are self-adjoint idempotents, $P_{+}P_{-} = 0$, commuting with $A_{1}$, $e_{1}$, $B_{1}$, $B_{2}$. Then $A_{2}'$ is self-adjoint, $A_{2}'^{2} = e_{1}^{2}P_{+} + p_{1}^{2}P_{-} = 1$, it commutes with $B_{1}, B_{2}$, and $A_{1}A_{2}' + A_{2}'A_{1} = (p_{1}e_{1} + e_{1}p_{1})P_{+} + 2P_{-} = 1 - p_{3}$, while $[A_{1},A_{2}'] = 2p_{1}e_{1}P_{+}$. Hence $C^{2} = 4 - 2p_{1}e_{1}P_{+}\cdot 2B_{1}B_{2} = 4 + 4KP_{+}$. Put $\Pi = \tfrac14(1 + K)(1 + p_{3})$, a self-adjoint idempotent (product of commuting ones) with $\Pi P_{-} = 0$, $\Pi K = \Pi$, so $C^{2}\Pi = 8\Pi$ and $(1 - p_{3})\Pi = 0$; $\Pi$ commutes with $C$ because $K$ and $p_{3}$ commute with each term. $\varphi_{r}(x) = 4\langle\Pi(1 + rC)x\rangle_{0}$. On the range of $P_{+}$ the configuration is that of (i): $CP_{+} = C_{0}P_{+}$ with $C_{0}$ the element of (i), whose four blades and their products with $K$ are not scalars and do not contain $p_{3}$; so $4\langle\Pi\rangle_{0} = 1$, $4\langle\Pi C\rangle_{0} = 0$, giving $\varphi_{r}(1) = 1$, $\varphi_{r}(C) = 4r\langle\Pi C^{2}\rangle_{0} = 8r$, $\varphi_{r}(C^{2}) = 8 + 8r\cdot 4\langle\Pi C\rangle_{0} = 8$, and $\varphi_{r}((1 - p_{3})^{2}) = 0$. Positivity: with $y = \Pi x$, $\tfrac14\varphi_{r}(xx^{\dagger}) = \langle y^{\dagger}y\rangle_{0} + r\langle y^{\dagger}Cy\rangle_{0}$ and, by the Cauchy-Schwarz inequality for the positive form $\langle a^{\dagger}b\rangle_{0}$, $\langle y^{\dagger}Cy\rangle_{0}^{2} \le \langle y^{\dagger}y\rangle_{0}\langle y^{\dagger}C^{2}y\rangle_{0} = 8\langle y^{\dagger}y\rangle_{0}^{2}$; so $8r^{2} \le 1$ gives $\varphi_{r}(xx^{\dagger}) \ge 0$.
